The primary negative finding is that the original UniProt scale-up labels
confound review status with class. Its $0.921$ held-out AUROC and $0.926$
APD6 comparison are descriptive, not validated measures of AMP recognition
on review-status-matched data (Section~\ref{sec:labelconfound}). Second, the chain
graph network underperforms the CNN at every scale we tested; without
contact-map or structure-derived edges, message passing along the backbone
adds parameters and optimization difficulty without adding information the
convolutions lack. We keep the GNN in the paper as a negative architectural
result. The small pilot scored about 0.95 AUROC while the larger, cluster-aware
experiment scored 0.92; because both data scale and split method changed,
this contrast does not by itself quantify homology leakage.

Limitations are real. The negative class is ``not annotated antimicrobial in
UniProt,'' which is not the same as experimentally non-antimicrobial; some negatives may be unrecognized AMPs and any bias direction depends on
which subsets are affected. The APD6 sensitivity at a 5\% negative-set FPR is 0.601; without
family-level analysis, we cannot attribute this miss rate to a specific
AMP class. The annealed designs and screen hits carry
no experimental validation; their novelty scores are sequence-level, not
structural. DBAASP was unreachable programmatically during this work (its
full-dataset endpoint returns an empty body without an interactive session),
so the APD6 list served as the sole external benchmark; adding DBAASP's
activity-annotated records is the natural next validation. Finally, the
sandbox budget (2 CPU, 2 GB RAM) capped model size, epochs, and the training
subsample; a larger training run might alter the numbers, but cannot by itself
remove the label-source confound.

Verdict: the UniProt scale-up is compromised by label-source confounding
and must be redone before making a clean recognition claim. The independent
published-corpus comparison, negative protocol lessons, and the ranked
candidates remain reproducible in-silico results. Nothing here is a
therapeutic, a wet-lab discovery, or a statistically established benchmark
record.
